% =====================================================================
%  5 Benutzerschnittstellen  (Aufgabenstellung: „Beschreibung der GUIs.
%  Struktur und Verhalten der wichtigsten Dialoge inkl. Inhalt und Aufbau
%  (z. B. Skizzen, Mock-ups, Screenshots), der Dialogflüsse und der
%  Eingabevalidierung (z. B. Gültigkeitsregeln für Eingabefelder).“)
%  Mock-ups in docs/diagramme/06_Mockups/, Export nach docs/diagramme/export/.
% =====================================================================
\section{Benutzerschnittstellen}
\TODO{Kurze Hinführung; Gestaltungsprinzipien (selbsterklärend, intuitiv) mit Quelle}

\subsection{Dialoge}
\TODO{Je wichtigem Dialog: Zweck, Inhalt, Aufbau, Verhalten; Mock-up einfügen}
\begin{figure}[H]
  \caption{Mock-up des Hauptfensters}\label{abb:mockup-haupt}
  \centering
  \grafik{mockup_hauptfenster}
  \quelle{Eigene Darstellung.}
\end{figure}

\subsection{Dialogfluss}
\TODO{Navigation zwischen den Dialogen (z. B. als Zustands- oder Ablaufdiagramm)}
\begin{figure}[H]
  \caption{Dialogfluss}\label{abb:dialogfluss}
  \centering
  \grafik{dialogfluss}
  \quelle{Eigene Darstellung.}
\end{figure}

\subsection{Eingabevalidierung}
\begin{longtable}{@{}L{3.2cm}L{2.4cm}L{5.4cm}L{4.2cm}@{}}
  \caption{Gültigkeitsregeln für Eingabefelder}\label{tab:validierung}\\
  \toprule
  \textbf{Dialog/Feld} & \textbf{Typ} & \textbf{Gültigkeitsregel} & \textbf{Fehlermeldung} \\
  \midrule
  \endfirsthead
  \toprule
  \textbf{Dialog/Feld} & \textbf{Typ} & \textbf{Gültigkeitsregel} & \textbf{Fehlermeldung} \\
  \midrule
  \endhead
  \bottomrule
  \endfoot
  \TODO{} & \TODO{} & \TODO{} & \TODO{} \\
\end{longtable}
\quelle{Eigene Darstellung.}
